% Generated draft: requires mathematical and source review.
\documentclass[11pt]{article}
\usepackage[margin=25mm]{geometry}
\usepackage{fontspec}
\setmainfont{Latin Modern Roman}
\usepackage{amsmath,amssymb,mathtools,unicode-math}
\setmathfont{Latin Modern Math}
\usepackage{enumitem,needspace,xurl,hyperref}
\setlength{\parindent}{0pt}
\setlength{\parskip}{5pt}
\setlength{\emergencystretch}{3em}
\newcommand{\N}{\mathbb{N}}
\newcommand{\Z}{\mathbb{Z}}
\newcommand{\Q}{\mathbb{Q}}
\newcommand{\R}{\mathbb{R}}
\newcommand{\C}{\mathbb{C}}
\newcommand{\F}{\mathbb{F}}
\newcommand{\A}{\mathbb{A}}
\newcommand{\PP}{\mathbb{P}}
\newcommand{\E}{\mathbb{E}}
\newcommand{\eps}{\varepsilon}
\newcommand{\dd}{\,\mathrm{d}}
\begin{document}
\section*{UnsolvedMath ID: 20003317}
\textbf{Source problem code:} catalogue entry
\subsection{Definitions and mathematical statement}
This record is the middle of a sentence after Conjecture 33, not another conjecture. The supplied source assessment and neighboring fragments identify its scope: D. Kleinbock observed that for two diagonal rays in opposite Weyl chambers, the argument of Kleinbock--Margulis [108] can be adapted to produce the witnesses of Conjecture 33, with a full-Hausdorff-dimensional witness set. Reference [108] itself is described by the aid as a one-flow thickness/avoidance result, not a theorem asserting an arbitrary two-ray intersection. The missing words before and after the fragment are not promoted to a standalone problem.

The mathematical lemma developed in the aid is local dimension subtraction. If $E,B$ are subsets of a $d$-dimensional manifold and, in a nonempty open $U$,
\[
 \dim_H(E\cap U)=d,\qquad\dim_H(B\cap U)<d,
\]
then
\[
 \dim_H((E\setminus B)\cap U)=d.
\]
This follows from $E\cap U=((E\setminus B)\cap U)\cup(E\cap B\cap U)$ and the maximum rule for the Hausdorff dimension of a finite union. Thus if $E$ is thick, meaning full dimension $d$ in every nonempty open set, it stays thick after deleting a globally lower-dimensional $B$, or a $B$ which is strictly dimension-subcritical in each such open set. Equality $\dim_HB=d$ would not suffice.

For the supplied modern application, let $X_3=SL(3,\R)/SL(3,\Z)$, of dimension $8$, and let $A$ be the full diagonal group, with positive identity component $A^0$. Fix a finite or countable family of nontrivial rays $A_i^+=\{\exp(tH_i):t\geq0\}$ and their two-sided one-parameter subgroups $A_i$. Put
\[
 E=\{x:A_ix\text{ is bounded for each prescribed }i\},\qquad
 B=\{x:A^0x\text{ is bounded}\}.
\]
Bounds and compact sets may depend on $x$ and $i$; a single compact set for all $i$ is not required. The aid invokes the An--Guan--Kleinbock joint-bounded-orbit theorem for thickness of $E$, and the Einsiedler--Katok--Lindenstrauss dimension-$2$ theorem for $B$. Finite components do not change full-$A$ boundedness. Therefore $E\setminus B$ is thick: every prescribed ray orbit is bounded there, while the full diagonal orbit is unbounded. Taking two rays gives the witness conclusion in the neighboring Conjecture 33. This is a conditional deduction from these precisely stated modern inputs, not a claim that [108] alone proves their countable intersection.

The original fragment concerns its historical opposite-chamber argument. The broader modern deduction uses different sources, whose complete hypotheses and citations must be checked; the elementary subtraction lemma by itself constructs no common bounded-orbit set.
\subsection{Short English statement}
The entry continues a historical claim about opposite-Weyl-chamber rays and is not an independent problem. The supplied aid isolates a valid dimension argument: deleting a strictly lower-dimensional set from a locally full-dimensional set preserves thickness. Combined with separate modern joint-boundedness and full-diagonal dimension theorems, it yields the desired witnesses without attributing a two-ray theorem to a one-flow source.
\subsection{Sources}
\begin{itemize}
\item Local UnsolvedMath registry, permanent ID $20003317$, source code catalogue entry; canonical location \texttt{aim-other-notes.json}, \texttt{notes[123]}. Neighboring records $20003315,20003316,20003318$ provide Conjecture 33 and the surrounding sentence boundaries.
\item AIM \emph{Emerging applications of measure rigidity}, continuation after Conjecture 33, \url{https://aimath.org/WWN/measrigid/measrigid.pdf}. It credits D. Kleinbock's adaptation of the argument in [108], identified in the aid as Kleinbock--Margulis, without a full bibliographic entry.
\item An--Guan--Kleinbock's 2015 common-bounded-orbit theorem and Einsiedler--Katok--Lindenstrauss, Theorem 10.2, for the dimension-$2$ full-diagonal bounded locus, as named by the supplied aid and neighboring records. Full publication details are absent.
\item Research attempt 1 is marked context-only/not-a-problem, proved, and low novelty confidence, despite the registry conjecture/open labels. Its explicit caveat says the aid is machine generated, not peer reviewed, and requires independent verification before citation.
\end{itemize}
\end{document}
